\chapter{Общий носитель гауссовой ковариации для физического повторного открытия}
\label{ch:v9-physical-reopening-common-origin-carrier}

\section{Один оператор для двух provenance-пакетов}

Пусть $R_\theta,R_\kappa>0$ --- два type-operator на пространстве
$V_{\rm type}\cong\mathbb R^5$ с multiplicities $1+1+3$ и
\begin{equation}
 \Tr R_\theta=\Tr R_\kappa=5.
 \label{eq:v9-common-origin-shape-normalization}
\end{equation}
На общем носителе $V_{\rm type}\otimes\mathbb R^2_{\rm KMS}$ введём
\begin{equation}
 Q(e,\chi,R_\theta,R_\kappa)
 =\operatorname{diag}(eR_\theta,e\chi^2R_\kappa),
 \qquad e=\frac{E_*}{\mu}>0.
 \label{eq:v9-common-origin-operator}
\end{equation}
Его dimension равна $10$, а два package-projector имеют rank $5$.

Один и тот же оператор содержит оба требуемых пакета:
\begin{equation}
 \Tr Q=5e(1+\chi^2),
 \qquad
 \det Q=e^{10}\chi^{10}\det R_\theta\det R_\kappa.
 \label{eq:v9-common-origin-trace-determinant}
\end{equation}

\section{Spectral-entropy parent}

Рассмотрим логарифм определителя functional
\begin{equation}
 \Phi(Q)=\Tr Q-\log\det Q-10.
 \label{eq:v9-common-origin-spectral-entropy}
\end{equation}
Для positive $Q$ это удвоенная относительная энтропия centered Gaussian
covariance $Q$ относительно unit covariance. Поэтому
$\Phi(Q)\geq0$ и equality достигается только при $Q=I_{10}$; связь
trace--логарифм определителя является стандартной Gaussian covariance identity
\cite{v9-common-origin-lami2017}.

В log-ratio coordinates $(x_1,x_2)$ и $(y_1,y_2)$ для двух shapes
стационарная точка равна
\begin{equation}
 (e,\chi,x_1,x_2,y_1,y_2)=(1,1,0,0,0,0).
 \label{eq:v9-common-origin-stationary-point}
\end{equation}
Точный Hessian имеет
\begin{equation}
 \rank\Hess\Phi=6,
 \qquad \det\Hess\Phi=36,
 \label{eq:v9-common-origin-hessian-rank}
\end{equation}
и spectrum
\begin{equation}
 \Spec(\Hess\Phi)=
 \left\{15-5\sqrt5,15+5\sqrt5,1,1,\frac35,\frac35\right\}.
 \label{eq:v9-common-origin-hessian-spectrum}
\end{equation}
Следовательно, один ограниченный снизу functional условно выбирает scale,
coupling и четыре relative-shape coordinates.

\section{Почему physical scale ещё не выведен}

Оператор зависит не от $E_*$ отдельно, а от отношения $e=E_*/\mu$.
Общее rescaling
\begin{equation}
 (E_*,\mu)\longmapsto(sE_*,s\mu)
 \qquad\Longrightarrow\qquad
 \frac{sE_*}{s\mu}=\frac{E_*}{\mu}
 \label{eq:v9-common-origin-scale-orbit}
\end{equation}
не меняет $Q$ или $\Phi$. Линеаризованная scale-map равна
\begin{equation}
 J_{\rm scale}=\begin{pmatrix}1&-1\end{pmatrix},
 \qquad \rank J_{\rm scale}=1,
 \qquad \operatorname{nullity}J_{\rm scale}=1.
 \label{eq:v9-common-origin-scale-nullity}
\end{equation}
Таким образом minimum $e=1$ означает лишь $E_*=\mu$. Без физического
происхождения reference covariance и её масштаба $\mu$ абсолютная энергия
не предсказана.

Кроме того, Gaussian covariance carrier пока является admitted auxiliary
carrier, а не состоянием, выведенным из существующего endpoint QMS. Поэтому
relative-entropy representation не равна physical parent-origin.

\begin{theorem}[Common-origin carrier admission]
Десятимерный positive covariance operator $Q$ является минимальным общим
carrier для scale/coupling и логарифм определителя packages. Functional $\Phi$ имеет
изолированный strictly stable minimum и условно закрывает joint selection
$2/2$.
\end{theorem}

\begin{nogocriterion}[Reference-state boundary]
Нельзя считать $E_*$ выведенным из minimum $e=1$, пока independently не
выведен опорный масштаб $\mu$. Нельзя также считать Gaussian relative
entropy физическим родительским функционалом, пока unit reference covariance не получена как
стационарное состояние того же microscopic carrier.
\end{nogocriterion}

\begin{protocol}\begin{enumerate}
\item common covariance carrier dimension: \textbf{$10$};
\item common carrier architecture: \textbf{$1/1$};
\item conditional joint selection: \textbf{$2/2$};
\item common Hessian rank/determinant: \textbf{$6/36$};
\item reference-scale nullity: \textbf{$1$};
\item physical-origin packages: \textbf{$0/2$};
\item physical reopening: \textbf{нет}.
\end{enumerate}\end{protocol}

Машинный сертификат сохранён в
\path{s2t_v9_physical_reopening_common_origin_carrier_admission_gate_results.json}.

\begin{thebibliography}{9}
\bibitem{v9-common-origin-lami2017}
L.~Lami, C.~Hirche, G.~Adesso, A.~Winter,
\emph{From log-determinant inequalities to Gaussian entanglement via
recoverability theory}, IEEE Transactions on Information Theory 63 (2017),
7553--7568; arXiv:1703.06149.
\end{thebibliography}